\pdfbookmark[0]{Prefacio}{Prefacio}

\chapter*{\textsf{\textbf{Prefacio}}}

En la última década, el estudio del microbioma humano ha pasado de la mera descripción de “quién está ahí” a interrogar “qué hace” y “cómo se relaciona con la salud”. La generalización de la secuenciación y la disponibilidad de cohortes públicas han permitido caracterizar comunidades microbianas en distintos nichos y asociarlas con fenotipos clínicos. Este avance ha venido acompañado de retos metodológicos bien conocidos: datos composicionales y ruidosos, alta variabilidad interindividual, heterogeneidad técnica entre estudios y la dificultad de separar asociación de causalidad. Aun así, el campo ha estudiado la idea de que ciertas alteraciones comunitarias pueden informar sobre el estado de salud y el riesgo de enfermedad, abriendo la puerta al desarrollo de biomarcadores microbianos con utilidad clínica. 

Esta Tesis Doctoral nace de una pregunta sencilla de formular y exigente de responder: ¿Podemos, a partir del microbioma de un nicho concreto, identificar los taxones relevantes y, con ello, ayudar a diagnosticar a un paciente? El reto es doble. Por un lado, de partida se desconoce qué taxones importan; por otro, al perfilar comunidades complejas, la señal biológica queda a menudo enmascarada por variabilidad inter e intrasujeto y ruido técnico. Este escenario demanda enfoques analíticos capaces de distinguir patrones robustos donde las señales son débiles o están correlacionadas.  

En cuanto a los motivos, esta investigación se asienta en una doble oportunidad. Por un lado, la madurez del ecosistema de datos públicos y de herramientas bioinformáticas permite desarrollar, reutilizar y comparar estudios con mayor rigor. Por otro lado, la clínica demanda criterios objetivos para clasificar, estratificar y anticipar respuestas terapéuticas. El microbioma puede contribuir a ambas necesidades. 

Con ese trasfondo, el objetivo central de la Tesis Doctoral fue explorar el microbioma humano como fuente de biomarcadores clínicos mediante metodologías bioinformáticas y de \textit{Machine Learning}, evaluando su capacidad para clasificar estados patológicos y apoyar la estratificación de pacientes. Este propósito se tradujo en objetivos específicos: (i) identificar patrones microbianos asociados a estados clínicos en distintos nichos anatómicos; (ii) construir modelos predictivos e interpretables que, además de clasificar, identifiquen las variables de mayor relevancia biológica; y (iii) evaluar la robustez y generalización de los hallazgos en marcos analíticos alternativos y, cuando es posible, en cohortes externas. 

Bajo este enfoque, la Tesis Doctoral explora tres escenarios complementarios. En enfermedades metabólicas (síndrome metabólico y diabetes tipo 2), el análisis pone de relieve señales microbianas diferenciadoras con potencial para orientar la selección de biomarcadores. En diabetes tipo 1 infantil, se construyen modelos predictivos interpretables en los que un conjunto acotado de taxones con alta relevancia biológica sustenta la decisión. Y en vaginosis bacteriana, se propone un subtipado alternativo del microbioma vaginal que aporta información complementaria a los métodos de estratificación existentes y puede apoyar, con cautela, la estimación de respuesta terapéutica. 

La contribución de esta Tesis Doctoral es tanto conceptual como práctica. Conceptual, porque muestra que el microbioma, aun sin asumir causalidad, puede proporcionar señales útiles para el diagnóstico y la estratificación. Práctica, porque ofrece flujos de análisis reproducibles e interpretables y consolida criterios de evaluación para distinguir hallazgos de posibles artefactos. Se discuten explícitamente alcance y limitaciones: la resolución taxonómica condicionada por la tecnología de perfilado, la naturaleza composicional de los datos y la heterogeneidad entre estudios. Estas consideraciones informan las recomendaciones finales para aproximar las firmas microbianas a escenarios clínicos realistas.